\documentclass[11pt,a4paper,sans]{moderncv}
%---------%---------%---------%---------%---------%---------%---------%---------
% Includes of packages, etc
%---------%---------%---------%---------%---------%---------%---------%---------

\usepackage[scaled]{helvet}
\renewcommand{\familydefault}{\sfdefault}
\moderncvstyle{banking}
\setlength{\hintscolumnwidth}{3cm}
\moderncvcolor{red}
\nopagenumbers{}

\usepackage[scale=0.75]{geometry}
\setlength{\footskip}{136.00005pt}

\usepackage{graphicx}
\graphicspath{{../includes}} % Specify the directory here
\usepackage{enumitem}

\ifxetexorluatex{}
  \usepackage{fontspec}
  \usepackage{unicode-math}
  \defaultfontfeatures{Ligatures=TeX}
  \setmainfont{Latin Modern Roman}
  \setsansfont{Latin Modern Sans}
  \setmonofont{Latin Modern Mono}
  \setmathfont{Latin Modern Math}
\else
  \usepackage[utf8]{inputenc}
  \usepackage[T1]{fontenc}
  \usepackage{lmodern}
\fi

\usepackage[swedish]{babel}
\usepackage{underscore}
% ------- % ------- % ------- % ------- % ------- % ------- % -------
% Self-defined commands
% ------- % ------- % ------- % ------- % ------- % ------- % -------

\newcommand{\thepage}{\arabic{page}}
\newlist{quotelist}{enumerate}{1}
\setlist[quotelist]{label=\textquotedblleft\arabic*\textquotedblright}

\makeatletter

% Definiera bredden om den inte finns
\@ifundefined{@photowidth}{%
  \newlength{\@photowidth}
  \setlength{\@photowidth}{4cm}
}{}

% Ange bildfil
\def\@photo{../../__2._photos/rickardaberg-service-tekniker.jpeg}

% Modifiera \makecvtitle om stilen är 'banking'
\@ifpackageloaded{moderncvstylebanking}{%
  \let\oldmakecvtitle\makecvtitle
  \renewcommand*{\makecvtitle}{%
    \vspace*{-1em}
    \begin{center}
      \includegraphics[width=\@photowidth,keepaspectratio]{\@photo}
    \end{center}
    \vspace{0.5em}
    \oldmakecvtitle
  }%
}

\makeatother

\newcommand{\mybeforetitle}{
  \vspace*{-2em}
  \begin{center}
    {\color{black}\Huge\sffamily\itshape \textbf{K}yl- och \textbf{V}ärmepumpmontör}
  \end{center}
  \begin{center}
    {\large\sffamily Utbildningssökande -- Personal Skills, Helsingborg \textbar{} Utförlig version}
  \end{center}
  \vspace{1em}
}

%------------------%------------------%------------------%------------------%------------------
% Title - photo etc
%------------------%------------------%------------------%------------------%------------------
\title{Civilingenjör Datateknik i karriärväxling mot kyl- och värmepumpteknik}
%------------------%------------------%------------------%------------------%------------------
% Address - email - phone
%------------------%------------------%------------------%------------------%------------------
\firstname{Rickard}
\lastname{Åberg}
\address{Bodekullsgatan 34b}{21440 Malmö}
\phone[mobile]{+46~709~43~14~01}
\email{raberg@duck.com}
%------------------%------------------%------------------%------------------%------------------
\quote{Fysiskt, tekniskt och konkret arbete -- där jag ser resultatet av dagens insats direkt.}
%------------------%------------------%------------------%------------------%------------------

\begin{document}
\mybeforetitle
\makecvtitle{}

\section{Sammanfattning}
Civilingenjör i datateknik (Linköpings tekniska högskola) med över 20 års erfarenhet av tekniskt arbete, systemfelsökning och problemlösning inom IT, mjukvaruutveckling och drift. Söker nu en konkret karriärväxling mot ett praktiskt hantverksyrke -- kyl- och värmepumpteknik -- där teknisk förståelse, noggrannhet och kundnära arbete kombineras med fysiska, handgripliga arbetsuppgifter och ett tydligt resultat varje dag. Van vid att snabbt sätta mig in i nya tekniska system, arbeta metodiskt enligt rutiner, checklistor och säkerhetskrav, samt dokumentera och felsöka avvikelser. B-körkort finns.

\section{Varför kyl- och värmepumpteknik}
\cvitem{}{
Efter många år i mer teoretiska IT- och konsultroller vill jag arbeta med något konkret och tekniskt tillsammans med händerna, där jag kan lämna arbetsdagen bakom mig när jobbet är klart. Kyl- och värmepumpteknik är ett hantverksyrke i stark efterfrågan, med tydlig koppling till klimatomställningen, och passar väl med mitt tekniska intresse för system, styrlogik, elektronik och mätning/kontroll. Jag deltog på Arbetsförmedlingens utbildningsmässa i Malmö den 9 september och gjorde ett studiebesök hos Personal Skills i Helsingborg för att lära mig mer om utbildningen till kyl- och värmepumpmontör.
}

\section{Teknisk bakgrund och systemförståelse}
\cvitem{}{
\begin{itemize}
  \item \textbf{Teknisk systemförståelse} -- Civilingenjör i datateknik med bred förståelse för hur tekniska system samverkar, inklusive mjukvara, hårdvara, sensorer, styrsystem och reglerteknik.
  \item \textbf{System- och driftorienterat arbetssätt} -- Van vid att arbeta nära tekniska system i drift, med fokus på stabilitet, tillgänglighet och korrekt funktion -- direkt överförbart till kylanläggningar och värmepumpar i drift.
  \item \textbf{Logiskt och strukturerat tänkande} -- Stark förmåga att analysera systembeteenden, följa flöden och förstå sekvenser, beroenden och styrlogik (relevant för kyl-/värmepumpstyrning, termostater och reglersystem).
  \item \textbf{Teknisk felsökning och avvikelsehantering} -- Metodisk i identifiering av felorsaker, bedömning av risker och genomförande av åtgärder vid tekniska avvikelser eller driftstörningar.
  \item \textbf{Elektronik och elteknik} -- Grundläggande erfarenhet av elektronik, mätning och praktisk felsökning sedan tidiga arbetslivserfarenheter inom bilelektronik samt hårdvarukurser under civilingenjörsutbildningen.
  \item \textbf{Mätning, noggrannhet och kontroll} -- Förstår vikten av precision, verifiering och uppföljning i tekniska och säkerhetskritiska sammanhang -- centralt vid köldmediehantering och täthetskontroller.
  \item \textbf{Rutiner och säkerhet} -- Van vid att arbeta enligt givna rutiner, checklistor och instruktioner, med högt säkerhetsmedvetande.
  \item \textbf{Dokumentation och rapportering} -- Erfaren av teknisk dokumentation, tydlig rapportering och spårbarhet av åtgärder och beslut.
  \item \textbf{Snabb inlärning i nya system} -- Snabblärd i nya tekniska system och arbetsmetoder, med fokus på korrekt hantering, kvalitet och säker drift.
\end{itemize}
}

\section{Praktiskt arbete och arbetsförmåga}
\cvitem{}{
\begin{itemize}[leftmargin=1.5em]
  \item Nyfiken på att lära sig nytt och med god studievana
  \item Vana vid fysiskt krävande arbete under längre arbetspass, bland annat som fältarbetare inom skogsbruk
  \item Arbetar metodiskt med fokus på säkerhet, ordning och kvalitet
  \item Självständig, uthållig och bekväm med praktiska, repetitiva arbetsmoment
  \item God förståelse för arbetsmiljö, riskbedömning och rutiner
  \item Trivs med arbete utomhus, hos kund och med tydligt avgränsade arbetsuppgifter
  \item God servicevana och van att förklara tekniska frågor för icke-tekniska kunder på ett tydligt sätt
\end{itemize}
}

\section{Yrkeserfarenhet}
\cventry{2026 --}{iOS-app: Breathwork, HRV \& Biofeedback}{Eget projekt}{}{}{
Design och utveckling av en iOS-app som integrerar HRV-analys (hjärtfrekvensvariabilitet), andningsstyrningsalgoritmer och realtidsbiofeedback. Visar fortsatt tekniskt driv och förmåga att självständigt driva ett projekt från idé till fungerande produkt.
}
\cventry{2024 --}{AI-verktyg och Claude}{Eget lärande}{}{}{
Praktisk, löpande erfarenhet av att arbeta med AI-assistenter (Claude/LLM-baserade verktyg) för utveckling, dokumentation, felsökning och automatisering. Exempel på snabb inlärning och anpassning till ny teknik -- samma förmåga som krävs för att sätta sig in i nya tekniska system inom kyl- och värmepumpbranschen.
}
\cventry{2026}{Fältarbetare (visstid)}{Skogsstyrelsen}{Sverige}{}{
Naturvårdsarbete, fältinventering, miljöövervakning och regelefterlevnad inom skogsbruk -- fysiskt utomhusarbete med hög grad av självständighet och rutinföljsamhet.
}
\cventry{2025}{Cloud Computing och Azure-kurs}{ReDI School of Digital Integration}{Malmö}{}{
Praktisk, kursfokuserad utbildning inom Azure och molnsäkerhet. Visar förmåga att snabbt tillgodogöra sig och certifiera sig inom ett helt nytt tekniskt område.
}
\cventry{2021--2022}{Fullstack-utvecklare}{ADBSafegate AB}{Malmö}{}{Angular, TypeScript, Azure, Scaled Agile Framework (SAFe)}
\cventry{2021}{Fullstack-utvecklare}{Zaplox AB}{Lund}{}{Vue.js, Node.js, JavaScript, CSS, AWS}
\cventry{2017--2018}{Senior mjukvaruutvecklare}{Verisure Innovation}{Malmö}{}{Java, Android, Docker/Kubernetes, drift och felsökning i produktionsmiljö}
\cventry{2012--2017}{Managing Consultant}{Bouvet Syd AB}{Malmö}{}{Java-utvecklare, teamledare, DevOps-ansvarig -- ansvar för leverans, kvalitet och driftsäkerhet}
\cventry{2004--2012}{Konsult}{Cybercom}{Malmö}{}{Java, systemintegration, test -- teknisk felsökning i komplexa driftmiljöer}
\cventry{2003--2004}{Forskningsingenjör}{École des Mines}{Nantes, Frankrike}{}{Hårdvarunära utveckling och systemintegration}
\cventry{2002}{Embedded-utvecklare}{HemoCue}{Ängelholm}{}{Utveckling nära hårdvara och mätteknik i medicintekniska instrument}

\section{Tidiga arbetslivserfarenheter och praktikplatser}
\cventry{1989--1994}{Tidningsutdelare}{Kvällsposten}{Skepparkroken och Skälderviken}{}{
Utdelning av morgontidningar i två distrikt nästan varje helg under fem års tid, från 12 års ålder. Ansvarstagande, punktlighet och logistik i alla väder.
}
\cventry{1991}{Praktik: verkstadsstöd}{Bjäresläp och Bilelektro AB}{Ängelholm}{}{
Verkstadsstöd inom släpvagnsservice och bilelektronik. Grundläggande el- och teknisk förståelse samt praktisk problemlösning -- tidig kontaktyta med hantverksyrken.
}
\cventry{1993--1997}{Ferie- och sommararbete}{Barkåkra kyrka}{Ängelholm}{}{
Vaktmästeri, trädgårdsarbete, målning, handgrävning och skötsel av kyrkogård under flera somrar -- praktiskt utomhusarbete med eget ansvar.
}
\cventry{1995}{Praktik: IT-support}{Datafirma}{Båstad}{}{
IT-support, installationer och teknisk kundservice.
}

\section{Utbildning}
\cventry{1996--2003}{Civilingenjör, Datateknik (300 hp)}{Linköpings tekniska högskola}{Linköping}{}{
Hårdvarukonstruktion, inbyggda system med mikrokontroller, avancerad datorarkitektur, realtidsprogrammering.
}
\cventry{2025}{Cloud Computing och Azure-säkerhet}{ReDI School of Digital Integration}{Malmö}{}{Microsoft Azure Fundamentals (AZ-900)}
\cventry{2023--2024}{Google Cybersecurity Fundamentals}{Google/Coursera}{Online}{}{}
\cventry{2016}{Programmering och problemlösning i Python (2.5 hp)}{Blekinge tekniska högskola}{Ronneby}{}{}
\cventry{1993--1996}{Naturvetenskapligt program}{Rönneskolan}{Ängelholm}{}{
Tekniskt tillval: programmering, elektronik, mikrodatorer, PC-support. Matematik upp till E, kemi A och B, fysik A och B.
}

\section{Körkort och certifikat}
\cventry{}{Körkort B}{Personbil}{}{}{}

\section{Språk}
\cvitemwithcomment{Svenska}{Modersmål}{}
\cvitemwithcomment{Engelska}{Mycket god}{}
\cvitemwithcomment{Franska}{God}{}

\section{Referenser}
Lämnas på begäran.

\end{document}
